\section{The compensated inversion identity used above}

For completeness, the Fourier identity in the fixed-target calculation
has the same short proof as in the manuscript. Let $\mu$ be a finite
signed measure with zero mass and first moment, and finite absolute
second moment. Then, for every real $H$,
\[
 \int(x-H)|x-H|\,d\mu(x)
 =-\frac4\pi\int_0^\infty
       \frac{\Im(e^{-itH}\widehat\mu(t))}{t^3}\,dt.
\]
The right side is absolutely convergent. Indeed, scaling and one
integration by parts give
\[
 \int_0^\infty\frac{ty-\sin(ty)}{t^3}\,dt
       =\frac\pi4y|y|.
\]
To check the constant, use $t^{-2}=\int_0^\infty s e^{-st}\,ds$ and
Tonelli to obtain
\[
 \int_0^\infty\frac{1-\cos t}{t^2}\,dt
 =\int_0^\infty\left(1-\frac{s^2}{s^2+1}\right)ds=\frac\pi2.
\]
For fixed $y$, the compensated integrand has its sign, so its absolute
integral is $\pi y^2/4$. Hence Fubini against $|\mu|$ is justified by
$\int(x-H)^2\,d|\mu|(x)<\infty$. The linear term vanishes by the
two zero moments, proving the identity and absolute convergence.

Applied to the sum law minus its Gaussian and cubic signed comparator,
this proves the inversion step in the conservation calculation. The
characteristic function of the latter comparator is
$e^{-nt^2/2}(1-in\gamma t^3/6)$; its mass and first moment agree
with those of the sum. This identity does not provide any numerical
upper bound on its integral. Such a bound still needs the stated
low-frequency and high-frequency estimates.
